\documentclass[colorinlistoftodos]{article}
\usepackage[utf8]{inputenc}
\usepackage{natbib}
\usepackage[spanish]{babel}
\usepackage[numbib,nottoc]{tocbibind}
\usepackage{graphicx}
\usepackage{longtable}
\usepackage[hidelinks]{hyperref}
\usepackage{todonotes}


\title{
    {\large Universidad de Buenos Aires\vspace{20pt}}\\
    {\includegraphics[scale=0.4]{images/uba.eps}}\\
    {\large\textbf{Facultad de Ciencias Exactas y Naturales}}\vspace{16pt}\\
    {\small\textbf{Especialización en Exploración de Datos\\y Descubrimiento del Conocimiento}\vspace{16pt}}\\
    {\includegraphics[scale=0.5]{images/master.eps}}\\
    {Reglamentación del acceso al aborto:\\selección de rasgos para un análisis discursivo}
}
\author{Autora:\\Fernández Urquiza, Macarena}
\date{Fecha:\todo{Definir fecha}}

\begin{document}
\clearpage\maketitle
\thispagestyle{empty}

\newpage
\tableofcontents

\newpage

\begin{abstract}
\todo[color=blue!40]{Escribir abstract}Abstract.
\end{abstract}

\section{Introducción}

\subsection{Antecedentes}

% Temas a abordar
%% Impacto de la reglamentación del acceso al aborto
%% Importancia del acceso, por parte de los ciudadanos, a los datos
%%  de incumbencia gubernamental y de organización social
%% Posibilidad de hacer un análisis discursivo con mayor cantidad
%%  de datos --> FerCa
%% Relevancia de la buena selección de features para el análisis
%%  de texto y, en este caso, del discurso
%% Paper seleccionado

\subsection{Objetivos}

% objetivos grales: podrían ser los mismos que los
% del paper (?)
% objetivos particulares: podría ser contrastar si
% las observaciones realizada en el paper se
% verifican para el español

% GRAL
% Realizar una selección de rasgos que permita
% capturar las particularidades discursivas de los
% partidos a favor de la despenalización del aborto
% y de los partidos en contra

% PARTICULAR
% - Cotejar si las observaciones detalladas en el paper
% se evidencian también para el caso del español
% - Los rasgos seleccionados permiten entrenar un
% modelo que, dado un discurso, distinga a qué
% tipo de partido pertenece su emisor

\section{Datos}
\subsection{Preprocesamiento}
Para este trabajo se utilizó la versión taquigráfica de la sesión número
23 (reunión 28) celebrada por la Cámara de Senadores, en la que se abordó la
regulación del acceso a la interrupción voluntaria del embarazo y a sus
posteriores cuidados. Los datos fueron obtenidos de la página del Senado de la Nación
Argentina\footnote{https://www.senado.gob.ar/parlamentario/sesiones/29-12-2020/28/}
utilizando el método de \textit{scrapping}.\par
En primer lugar, se descargó la transcripción en formato \textit{PDF} y, luego, por
medio de la librería \textit{pdfminer}\footnote{https://pdfminersix.readthedocs.io/en/latest/},
se lo convirtió a texto plano, de manera que fuese procesable. Dado que la librería
empleada en la conversión no arrojó una conversión limpia y que, además, no todo el texto
resultaba de relevancia para el presente trabajo, se realizó una tarea de
preprocesamiento, con el fin de limpiar y organizar los datos de forma útil a los
objetivos aquí perseguidos.\par
Como primer paso del preprocesamiento se quitaron los encabezados y pie de paǵina.
El archivo \textit{.pdf}, convertido a \textit{.txt}, no contaba con separación de
páginas, por lo que los encabezados y pies se convertían en líneas de texto que
interrumpían los discursos de los senadores y participantes del debate.
Posteriormente, se extrajo la sección del texto pertinente para este trabajo:
la sección 6, dedicada a la ``Regulación  del  acceso  a  la  interrupción
voluntaria  del  embarazo  y  a  la atención postaborto''. Otras secciones consistían
en el izamineto de la bandera, la convocatoria a la sesión, la lectura de la ley
resultante, entre otras cuestiones que no hacían al discurso argumentativos de los
asistentes, sino más bien a cuestiones protocolares, por lo que fueron desestimadas
para este análisis.\par
Con el texto de interés delimitado, y mediante el uso de patrones regulares, se
identificó a los distintos oradores y sus respectivos discursos, como así también
a las secciones que no pertenecían a fragmentos discursivos emitidos durante la
discusión, sino a comentarios agregados por el taquígrafo\footnote{``Luego de unos
instantes'' y ``Contenido no inteligible'' constituyen ejemplos de estos
fragmentos.}. Siguiendo la metodología descripta por \cite{monroe2008fightin},
los datos fueron guardados en un archivo en formato \textit{XML}, en el que
cada fragmento fue etiquetado con las siguiente información:
\begin{itemize}
    \item \texttt{speech}: etiqueta que indica si el fragmento corresponde a un
    discurso efectivamente emitido durante la discusión, en cuyo caso toma el
    valor \texttt{true}, o si consiste en un comentario del taquígrafo, ante lo cual
    toma el valor \texttt{false}.
    \item \texttt{speaker}: en el caso de que el contenido del texto refiera a
    un fragmento discursivo, esta etiqueta indica el nombre del orador que lo
    pronunció. Si el fragmento es un comentario del traquígrafo, la etiqueta
    recibe el valor \texttt{none}.
\end{itemize}
\par
Adicionalmente, como resultado del preprocesamiento, se extrajo la lista de
asistentes a la sesión en formato tabular, con el objetivo de poder
enriquecerla adicionando especificaciónes como la filiación partidaria y la
decisión de voto.\todo[color=blue!40]{Desarrollar y explicar el scrapping y mergeo de la data.}

% =================================================
% Breve descripción de los datos:
% - cantidad de partidos totales
% - cantidad de oradores totales
% - distribución oradores~partidos
% - distribución género~partidos
% - palabras y tokens totales
% - distribución de palabras y tokens por partido
% - distribución de palabras y tokens por orador
\subsection{Descripción}
La sesión analizada contó con un total de 70 senadores,
pertenecientes a 26 partidos políticos, distribuidos del
siguiente modo (Figura \ref{fig-distrib-senators}):

\begin{figure}[h!]
\centering
\includegraphics[scale=0.5]{../visualizations/senators_distribution.png}
\caption{Distribución de senadores en partidos políticos y en provincias.}
\label{fig-distrib-senators}
\end{figure}

Respecto de los partidos, solo uno (Frente de todos) cuenta con 15
senadores; también uno solo (Alianza frente para la victoria) es
representado por 10 senadores; dor partidos (Alianza cambiemos y
Juntos por el cambio) cuentan con 7 senadores, y el resto de los
partidos tienen entre 3, 2 y un senador. Y en cuanto a las provincias
(24 en total), la mayoría tiene 3 senadores exceptuando a Tucumán y
a La Rioja, que tienen solo 2.\todo[color=pink]{Agregar tablas en anexo
dependiendo de la extensión del trabajo.}

La figura \ref{fig-distrib-senators} nos muestra además que la intención
de voto no guarda una relación unívoca con los partidos a los cuales los
senadores representan. A excepción de los partidos que cuentan con un único
senador, la mayoría cuenta con senadores que votaron a favor y senadores
que votaron en contra de la ley para el acceso al aborto.

\begin{figure}[h!]
\centering
\includegraphics[scale=0.5]{../visualizations/senators_vote.png}
\caption{Distribución votos en los distintos partidos políticos.}
\label{fig-distrib-vote}
\end{figure}

Por último, al considerar las intervenciones discursivas, observamos que
9 senadores no intervinieron en la sesión, y las intervenciones de los
restantes 61 senadores fueron resumidas en la Tabla \ref{table-tokens}, que plasma
las métricas univariadas de la cantidad de tokens totales y únicos y nos
permite observar que ambas obedecen una distribución sesgada a izquierda.

\begin{figure}[h!]
\centering
\includegraphics[scale=0.6]{../visualizations/tokens_vote.png}
\caption{Distribución de tokens emitidos por los oradores en relación a
su intención de voto.}
\label{fig-distrib-tokens-vote}
\end{figure}

La Figura \ref{fig-distrib-tokens-vote}, por otro lado, nos muestra,
de forma más detallada, la distribución de los tokens emitidos según
la intención de voto de los oradores.

\begin{table}[h!]
\begin{center}
\begin{tabular}{ |c|c|c|c| }
\hline
Tokens & Media & Mediana & Desvío Estándar \\
\hline\hline
Totales & 811.65 & 636.12 & 714.79 \\
\hline
Únicos & 296.47 & 254.25 & 238.68 \\
\hline
\end{tabular}
\caption{Metrícas univariadas de la cantidad de tokens}
\label{table-tokens}
\end{center}
\end{table}


\section{Metodología}
\subsection{Clasificación}
\subsection{Métodos estadísticos}
\subsection{Métodos predictivos}
% Describir métodos aplicados.

%% Clasificación ==> w1: 0.004, w2: 0.002
%   - RF
%% Nonmodel-based ==> d(D-R) = w>0=D | w<0=R
%   - Dif. de frecuencias absolutas
%   - Dif. de proporciones
%   - Dif. de proporciones (s/stopwords)
%   - Odds
%   - Log-odds ratio
%   - Log-odds ratio (s/baja frecuencia)
%   - TF-IDF
%   - Word Scores
%% Model-based ???

%%% NaiveBayes (baseline)
%%% Reg. Logística

%       | RF | DFA | DFR |
% NB    |                |
% RL    |    |     |     |

\section{Resultados}

\section{Discusión y conclusiones}

\bibliographystyle{linquiry3}
\bibliography{bibliography}

\listoftodos%

\end{document}
